\section{Closed-Loop Propagation of VLA Attacks}
\label{sec:taxonomy}

\subsection{Closed-Loop VLA System}

A VLA-enabled robot should be read as a closed-loop state transition system, not as a single prompt-to-action function. Let $x_t$ denote the physical state, $o_t=O(x_t,\eta_t)$ the observation produced by sensors and preprocessing noise $\eta_t$, $m_t$ the memory/tool/human context, and $u_t$ the command emitted by the VLA or its planner. A low-level controller $K$ maps $u_t$ and robot state to actuation, while a recovery or safety mechanism $R$ may stop, replan, filter, or ask for help:
\begin{align}
    u_t &= \pi_\theta(o_{\le t}, g, m_t), \nonumber\\
    a_t &= K(u_t, x_t, R_t), \nonumber\\
    x_{t+1} &= F(x_t, a_t, w_t), \quad m_{t+1}=U(m_t,o_t,u_t,a_t).
    \label{eq:closed-loop-vla}
\end{align}
The safety envelope is a set $\mathcal{S}$ over physical state, controller state, tool authorization, privacy state, and human proximity. An attack is relevant when an intervention $\delta$ changes one of these transitions or contexts so that the rollout violates a task, authorization, privacy, or safety predicate. This view gives the survey its central question: where does the adversary enter, which representation changes first, how does that change cross the controller boundary, and does feedback or recovery stop it?

Figure~\ref{fig:vla-pipeline} is the organizing framework for the rest of the survey. It separates where the adversary enters from what the attack changes, marks the controller boundary, and summarizes the propagation chain
\begin{align}
\delta &\rightarrow \Delta o_t \rightarrow \Delta z_t
\rightarrow \Delta u_t \rightarrow \Delta a_t \nonumber\\
&\rightarrow \Delta x_{t+1}
\rightarrow \text{recovery or violation}
\label{eq:attack-propagation}
\end{align}
The same perturbation can stop at the observation, change an embedding or grounding representation, alter a plan, shift an action token, corrupt a chunk, redirect a trajectory, be clipped by a controller, or only become visible after the next observation. This is why a VLA attack claim is not settled at the model-output boundary; the robot-level claim depends on execution, feedback, recovery, and the resulting physical or system consequence.

\subsection{Attack Map}

The attack map has four intuitive stages. First, the attack enters through a channel: vision, language, memory, tool output, middleware, training data, physical environment, or human interaction. Second, it changes a target representation: grounding, plan, reasoning trace, action token, action chunk, waypoint, trajectory, velocity field, controller state, memory entry, or tool call. Third, the changed representation reaches a robot or system consequence: task failure, wrong-object action, unsafe motion, unauthorized skill use, privacy leakage, or persistent state compromise. Fourth, the validation setting determines what the claim can support: digital input, real-sensor input, simulation rollout, hardware-in-the-loop, real-robot actuation, or human-proximate study.

Table~\ref{tab:taxonomy-overview} keeps these stages orthogonal because common labels hide different mechanisms. ``Physical attack'' describes an access channel, not the target: a physical patch may attack visual embeddings, grounding, visual proprioception, or a trajectory objective. ``Jailbreak'' describes a family resemblance, not a consequence: a prompt attack may cause task failure, targeted motion, unauthorized tool use, or no robot-level effect. ``Backdoor'' describes timing and training access, but the robotic target may be a gripper primitive, an action chunk, a waypoint, or a long-horizon trajectory.

\begin{table*}[t]
\centering
\footnotesize
\caption{Orthogonal multi-axis taxonomy for coding VLA attacks. Each paper may receive multiple labels on each axis; axes are not mutually exclusive categories.}
\label{tab:taxonomy-overview}
\renewcommand{\arraystretch}{1.08}
\setlength{\tabcolsep}{2.4pt}
\begin{tabularx}{\textwidth}{@{}p{0.16\textwidth}Y Y@{}}
\toprule
\textbf{Axis} & \textbf{Allowed labels} & \textbf{Coding question} \\
\midrule
Lifecycle & Training; inference; runtime; post-deployment & When is the adversarial intervention introduced or activated? \\
Entry surface & Vision; language; state; memory; tool; middleware; human; environment & Which semantic or system channel carries the adversarial content? \\
Target component & Visual encoder; language encoder; multimodal projector; planner/reasoner; action head; controller; memory; tool/API; middleware & Which system component is attacked or measured? \\
Target representation & Embedding; grounding; plan; action token; chunk; waypoint; trajectory; velocity field; controller state; memory entry & Which representation or control object changes inside or across components? \\
Security property violated & Availability; integrity; confidentiality; authorization & Which security property is violated? Physical safety is represented as a consequence rather than merged into this axis. \\
Knowledge level & White-box; gray-box; black-box; transfer-only; unknown & What does the attacker know about weights, gradients, architecture, prompts, action heads, or controllers? \\
Access channel & Digital; physical; human-mediated; supply-chain; hybrid & How does the attacker deliver the intervention? This is separate from knowledge level. \\
Privilege level & User; bystander; operator; data contributor; model provider; tool/API; middleware; no privilege & What role or privilege boundary does the attacker occupy or abuse? \\
Temporal profile & One-shot; repeated; delayed; persistent & Does the attack act once, recur over feedback, activate later, or survive sessions/training? \\
Validation tier & Digital input; real sensor; simulation; HIL; robot actuation; human-proximate & What is the strongest demonstrated validation tier? Do not infer actuation or human proximity from input-only tests. \\
Consequence & Task failure; wrong object; unsafe action; collision/contact risk; privacy leakage; unauthorized skill; persistent state compromise & What happened at the robot/system boundary? Keep consequence separate from objective and entry surface. \\
\bottomrule
\end{tabularx}
\end{table*}


\subsection{Typed Attack Success}
\label{subsec:attack-success}

Because VLA safety papers do not yet use a uniform definition of attack success, a survey needs a typed definition. Let $\tau^0(c)$ be the benign rollout from an initial condition and task context $c=(x_0,g,m_0)$, and let $\tau^\delta(c)$ be the rollout under an adversarial intervention $\delta$ that satisfies the declared knowledge, access, privilege, budget, and physical constraints. A VLA attack is successful for objective type $O$ and consequence class $C$ when
\begin{align}
    \mathrm{Succ}_{O,C}(\tau^\delta,\tau^0)
    &= \mathbf{1}\{ \phi_{O,C}(\tau^\delta)=1 \nonumber\\
    &\quad \land\ \phi_{O,C}(\tau^0)=0
    \land\ \delta \in \Delta(\mathcal{A}) \}.
    \label{eq:attack-success}
\end{align}
where $\phi_{O,C}$ is a typed violation predicate and $\Delta(\mathcal{A})$ is the admissible attacker set. The benign rollout condition prevents ordinary task difficulty or non-adversarial distribution shift from being counted as attack success. If a benign rollout is unavailable, the paper should state the baseline used, such as clean task success, human-labeled intended behavior, or non-member trajectory statistics.

The predicate $\phi_{O,C}$ must be typed. A \emph{task-failure attack} succeeds when the benign task would complete but the attacked rollout fails. A \emph{targeted behavior attack} succeeds when the robot reaches an attacker-specified object, action primitive, waypoint, trajectory, skill, or tool call. A \emph{safety attack} succeeds when the rollout violates a safety envelope, contact margin, force/velocity/workspace bound, or human-proximity constraint. An \emph{authorization attack} succeeds when the system executes an unauthorized command, skill, tool, or privilege transition. A \emph{privacy attack} succeeds when membership, trajectory, user, or scene information is inferred above a stated threshold. A \emph{persistence attack} succeeds when the adversarial effect survives replanning, feedback, memory update, fine-tuning, or a session boundary. A \emph{recovery-failure attack} succeeds when the system has an opportunity to stop, replan, ask for help, or invoke a recovery policy but still reaches the violation predicate.

This definition makes ASR meaningful only when five items are stated together: the violation predicate, the denominator, the attacker assumptions, the validation tier, and the benign or clean baseline. Without these items, ASR values from different VLA papers should be treated as local evidence rather than comparable field-level rates.

\begin{figure*}[t]
\centering
\scriptsize
\begin{tikzpicture}[
    node distance=0.45cm and 0.55cm,
    box/.style={draw, rounded corners, align=center, minimum height=0.95cm, text width=2.15cm, fill=gray!8},
    atk/.style={draw, rounded corners, align=center, text width=2.2cm, fill=red!8},
    arr/.style={-{Latex[length=2mm]}, thick}
]
\node[box] (obs) {Entry point\\observation / source};
\node[atk, left=of obs, text width=1.9cm] (adv) {Adversary\\and attack\\operator $\delta$};
\node[box, right=of obs] (latent) {Changed\\representation\\$z_t$};
\node[box, right=of latent] (cmd) {Command / action\\$u_t$};
\node[box, right=of cmd] (ctrl) {Controller\\$K(u_t,x_t)$};
\node[box, right=of ctrl] (exec) {Physical state\\$x_{t+1}$};
\node[box, below=of cmd, text width=2.6cm] (rec) {Feedback and recovery\\memory / tools / human};
\node[box, below=of exec, text width=2.25cm, fill=green!8] (safe) {Recovery or\\violation\\$\mathcal{S}$};
\draw[arr] (obs) -- (latent);
\draw[arr] (latent) -- (cmd);
\draw[arr] (cmd) -- (ctrl);
\draw[arr] (ctrl) -- (exec);
\draw[arr] (exec.south) -- (safe.north);
\draw[arr] (exec.south) |- (rec.east);
\draw[arr] (rec.west) -| (latent.south);
\draw[arr] (rec.north) -- (cmd.south);
\draw[arr, red!70] (adv) -- node[above, sloped] {entry} (obs);
\draw[arr, red!70] (adv.north) .. controls +(0,1.3) and +(0,1.3) .. node[above] {language / tool / data / physical entries} (cmd.north);
\draw[arr, red!70] (adv.south) .. controls +(0,-1.1) and +(-1.3,-0.8) .. node[below] {middleware / recovery / memory entries} (rec.west);
\draw[arr, red!70, dashed] (latent.north) to[bend left=20] node[above] {$\Delta z_t$} (cmd.north);
\draw[arr, red!70, dashed] (cmd.north) to[bend left=20] node[above] {$\Delta u_t$} (ctrl.north);
\draw[arr, red!70, dashed] (ctrl.north) to[bend left=20] node[above] {$\Delta a_t$} (exec.north);
\end{tikzpicture}
\caption{Core closed-loop view used throughout the survey. The figure separates adversarial entry surfaces from representation impact, marks the controller boundary between model command and physical actuation, and shows how feedback either supports recovery or carries the system toward a robot/system violation.}
\label{fig:vla-pipeline}
\end{figure*}

\subsection{Attacker Assumptions}

We assume the deployer has an intended task, authorized users, authorized tools and skills, and a safety envelope for people, property, data, and robot hardware. The attacker seeks to violate one or more of these constraints while preserving enough surface plausibility that the system continues acting, logs the event as normal, or stores corrupted context for future sessions.

Attacker capability has three independent dimensions. \emph{Knowledge level} describes what the attacker knows about weights, gradients, architecture, prompts, action heads, or controllers: white-box, gray-box, black-box, transfer-only, or unknown. \emph{Access channel} describes how the intervention is delivered: digital, physical, human-mediated, supply-chain, or hybrid. \emph{Privilege level} describes the role or boundary involved: user, bystander, operator, data contributor, model provider, tool/API, middleware, or no privilege. Keeping these dimensions separate prevents a common confusion: ``black-box'' and ``physical'' are not alternatives. A physical patch can be optimized with white-box gradients, transferred in a black-box setting, or deployed by a bystander with no digital privileges.

\subsection{Reading the Map Through Examples}

The map is easiest to read through examples. In a scene-text hijacking scenario, the attacker does not need to rewrite the user's goal; the robot sees an instruction-like phrase in the environment, the phrase changes source grounding or plan selection, and the downstream consequence may be wrong-object manipulation or unauthorized skill use. In a patch scenario, the entry is visual, but the relevant question is whether the perturbation merely changes an embedding or whether it reaches an action token, waypoint, or physical trajectory \cite{wang2025vlaadversarial,xu2025edpa,lu2026patch}. In a poisoned-demonstration scenario, the attack enters before deployment, hides behind clean utility, and later activates as a gripper primitive, action freeze, or chunk-level drift \cite{wang2025freezevla,xu2025dropvla,xu2026silentdrift}. In a trajectory-redirection scenario, the language may remain close to the user's instruction while the physical path bends toward a different endpoint \cite{puthumanaillam2026trajectory}. In a membership-inference scenario, the consequence is not motion failure at all; action likelihoods or temporal action traces leak whether a demonstration or trajectory was in the training set \cite{peng2026miavla,luan2026vlaleaks}.

This separation also prevents overclaiming. A VLM jailbreak with no robot action is contextual evidence for source confusion, not direct evidence of VLA unsafe action. A ROS2 middleware exploit is embodied-adjacent evidence for tool/middleware compromise, not a VLA policy attack unless the VLA system consumes or emits the compromised message. A physical LiDAR or camera attack can be strong sensor evidence but remains contextual for VLA action unless a VLA-controlled robot rollout is measured.
